% =====================================================================
%  MLE-AI-201 Machine Learning -- Final Project report (IEEE-style)
%  Every number in the text comes from tables/numbers.tex and the other
%  tables/*.tex files, which `python -m src.run_all` regenerates.
%  Build:   pdflatex report.tex ; pdflatex report.tex
% =====================================================================
\documentclass[10pt,twocolumn,a4paper]{article}

\usepackage[a4paper,top=0.75in,bottom=1.0in,left=0.68in,right=0.68in,columnsep=0.25in]{geometry}
\usepackage{mathptmx}
\usepackage[T1]{fontenc}
\usepackage[utf8]{inputenc}
\usepackage{graphicx}
\usepackage{amsmath,amssymb}
\usepackage{booktabs}
\usepackage{array}
\usepackage{titlesec}
\usepackage{caption}
\usepackage{cite}
\usepackage[hidelinks]{hyperref}
\usepackage{url}

\renewcommand{\thesection}{\Roman{section}}
\renewcommand{\thesubsection}{\Alph{subsection}}
\titleformat{\section}[block]{\centering\normalsize\scshape}{\thesection.}{0.5em}{}
\titleformat{\subsection}[block]{\normalsize\itshape}{\thesubsection.}{0.5em}{}
\titlespacing*{\section}{0pt}{1.4ex plus 1ex minus .2ex}{1ex}
\titlespacing*{\subsection}{0pt}{1.0ex plus 1ex minus .2ex}{0.6ex}
\newcommand{\abstractkw}[1]{\noindent\textbf{\textit{#1}}}
\captionsetup{font=small,labelsep=period}
\setlength{\parindent}{1em}
\graphicspath{{figures/}}
\newcommand{\schwa}{\raisebox{\depth}{\rotatebox{180}{e}}}   % Azerbaijani letter schwa

\newcommand{\NRaw}{100,775}
\newcommand{\NClean}{57,422}
\newcommand{\NFeat}{99}
\newcommand{\NTrain}{40,196}
\newcommand{\NVal}{8,613}
\newcommand{\NTest}{8,613}
\newcommand{\TierThr}{209,000}
\newcommand{\TreeADepth}{20}
\newcommand{\TreeALeaf}{5}
\newcommand{\TreeBCrit}{entropy}
\newcommand{\TreeBDepth}{20}
\newcommand{\TreeBLeaf}{1}
\newcommand{\SvmLam}{0.001}
\newcommand{\RffGamma}{0.003}
\newcommand{\RffLam}{1e-05}
\newcommand{\OursARmse}{0.231}
\newcommand{\SkARmse}{0.232}
\newcommand{\OursARsq}{0.869}
\newcommand{\OursAMape}{14.7}
\newcommand{\OursAMae}{46.1}
\newcommand{\BaseARmse}{0.636}
\newcommand{\OursBTreeFone}{0.923}
\newcommand{\SkBTreeFone}{0.927}
\newcommand{\OursSvmFone}{0.909}
\newcommand{\OursSvmAuc}{0.968}
\newcommand{\OursRffFone}{0.913}
\newcommand{\OursRffAuc}{0.974}
\newcommand{\SgdFone}{0.907}
\newcommand{\SvcLinFone}{0.904}
\newcommand{\SvcRbfFone}{0.910}
\newcommand{\OursTreeAFit}{17.3}
\newcommand{\SkTreeAFit}{0.70}
\newcommand{\OursSvmFit}{5.0}
\newcommand{\SvcRbfFit}{7.0}
\newcommand{\SvmNormW}{3.386}
\newcommand{\SvmMargin}{0.591}
\newcommand{\SvmSvFrac}{25.7}
\newcommand{\SvmErrNear}{77}
\newcommand{\DistHouseR}{0.85}
\newcommand{\DistSizeR}{-0.55}
\newcommand{\SameA}{93.6}
\newcommand{\SameB}{98.2}
\newcommand{\GiniBest}{0.9227}
\newcommand{\EntropyBest}{0.9246}


\begin{document}

\twocolumn[
  \begin{@twocolumnfalse}
  \begin{center}
    {\LARGE\bfseries Pricing Baku Homes From Scratch:\\A NumPy Decision Tree and Pegasos SVM on bina.az Listings\par}
    \vspace{1.2em}
    % TODO(team): authors, Surname, Given name, alphabetical by surname.
    {\large
      Surname1, Name1\quad Surname2, Name2\quad Surname3, Name3\par}
    \vspace{0.5em}
    {\itshape MLE-AI-201 Machine Learning, AI Academy --- Cohort II\par}
    \vspace{1.4em}
  \end{center}
  \begin{quote}
    \small
    \abstractkw{Abstract---}%
    We clean a scraped dump of \NRaw{} bina.az sale records into \NClean{} unique Baku
    residential listings. We then frame two tasks: regression of $\log(\text{price})$, and a
    premium/standard price-tier label split at the training-set median (\TierThr{}~AZN).
    Both models are implemented in NumPy. The first is a CART decision tree with Gini, entropy and MSE
    criteria, a vectorised threshold search and four stopping rules. The second is a soft-margin SVM trained
    by Pegasos sub-gradient descent on the hinge-loss primal, extended to an RBF kernel with
    random Fourier features. Hyperparameters are chosen on a validation split only. On the held-out test set
    our tree reaches RMSE \OursARmse{} on log price ($R^2=\OursARsq$, mean absolute
    percentage error \OursAMape\% in AZN). scikit-learn's tree with the same settings scores \SkARmse. On the tier task our tree reaches
    F1 \OursBTreeFone{} (sklearn \SkBTreeFone), our linear SVM \OursSvmFone{} (SGDClassifier
    \SgdFone) and our RFF-SVM \OursRffFone{} (kernel SVC \SvcRbfFone). The from-scratch models
    match the library baselines within noise. Most remaining errors are luxury properties,
    house-heavy suburban districts and listings priced close to the tier threshold.
    Code: \href{https://github.com/TODO-team/TODO-repo}{\texttt{github.com/TODO-team/TODO-repo}}
    (tag \texttt{v1.0-final}).

    \vspace{0.6em}
    \abstractkw{Index Terms---}decision tree, support-vector machine,
    Pegasos, random Fourier features, real-estate price prediction.
  \end{quote}
  \vspace{1.0em}
  \end{@twocolumnfalse}
]

% =====================================================================
\section{Introduction}
Online listing prices are noisy, heavy-tailed and dominated by location. Predicting them is a
realistic test of whether a model has learned the market or just memorised the data. Our contributions:
(i) a documented, leakage-free cleaning pipeline for the raw bina.az dump;
(ii) a from-scratch CART tree whose splits match scikit-learn's exactly on a fully determined test example;
(iii) a from-scratch Pegasos SVM~\cite{pegasos} with visible margins and support vectors, plus a
random-Fourier-feature (RFF) kernel extension~\cite{rff};
(iv) a matched benchmark against scikit-learn~\cite{sklearn} with cross-validated uncertainty,
timing and an error analysis.
\texttt{python -m src.run\_all} regenerates every number, table and figure in this report.

\section{Background}
A classification tree chooses, at each node, the split $x_j\le t$ that maximises the information gain
$\mathrm{IG}=I(\text{parent})-\sum_c \tfrac{N_c}{N} I(c)$, where $I$ is the Gini impurity $1-\sum_k p_k^2$
or the entropy $-\sum_k p_k\log_2 p_k$. A regression tree uses the variance
$I_{\text{reg}}=\tfrac1N\sum_i(y_i-\bar y)^2$ instead~\cite{cart}. The soft-margin linear SVM minimises
the regularised hinge loss
\begin{equation}
  J(\mathbf{w},b)=\tfrac{\lambda}{2}\lVert\mathbf{w}\rVert^{2}
  +\tfrac{1}{n}\sum_{i=1}^{n}\max\!\big(0,\,1-y_i(\mathbf{w}^{\top}\mathbf{x}_i+b)\big),
\end{equation}
which Pegasos solves by stochastic sub-gradient steps with rate $\eta_t=1/(\lambda t)$, with no QP or SMO
solver~\cite{pegasos}. This objective equals the textbook $C$-SVM with $C=1/(\lambda n)$.

\section{Data and Preprocessing}
\subsection{Dataset and cleaning}
The Kaggle dump (\texttt{house\_sale.csv}, version~1, scraped Oct--Nov 2024) has \NRaw{} rows and 51 columns.
Most columns are strings with Azerbaijani names and units. The cleaning steps, in order, are logged by \texttt{data\_prep.clean}:
\begin{itemize}\setlength\itemsep{0pt}
\item \emph{Currency:} every price is AZN (checked), so no conversion is needed.
\item \emph{Scope:} we keep residential categories (new build, old build, house/villa) in Baku.
  Land, offices, garages and commercial objects are dropped because they form different markets,
  and land is priced per \emph{sot}.
\item \emph{Duplicates:} the same listing is re-scraped on many days. We keep the newest scrape per URL,
  which removes 33,054 rows, and then drop 709 re-posts with identical price, area, rooms, floor,
  coordinates and category.
\item \emph{Type coercion:} \texttt{Sah\schwa{}} (``145 m\textsuperscript{2}'')$\to$ area;
  \texttt{Torpaq sah\schwa{}si} (``1.3 sot'')$\to$ land; \texttt{M\schwa{}rt\schwa{}b\schwa{}} (``7 / 9'')$\to$ floor and total floors;
  yes/no columns$\to$0/1. The district (\emph{rayon}) and metro proximity are parsed from \texttt{extra\_info}.
\item \emph{Outliers:} we use fixed domain rules, not data quantiles, so no statistic is learned from test rows.
  The rules drop rows with price outside $[10^4, 5{\cdot}10^6]$~AZN, area outside $[15,1500]$~m\textsuperscript{2},
  price per m\textsuperscript{2} outside $[150,15000]$ (typos or rent ads), land above 200~sot,
  more than 15 rooms, impossible floors, or coordinates outside Baku. Together they remove 133 rows.
\item \emph{Missing values:} median imputation using \emph{training} medians, with a missing-indicator
  column for rooms, floor and repair. Missing land area means a flat, so it is set to 0.
\end{itemize}

\subsection{Leakage}
\texttt{unit\_price} (AZN per m\textsuperscript{2}) is price divided by area, and \texttt{total\_price}
equals \texttt{price} in 99.99\% of rows. Both are arithmetic functions of the target and are dropped. Including either one would
let the tree reconstruct the price exactly. We also drop identifiers and personal fields
(\texttt{owner\_name}, \texttt{shop\_name}, \texttt{address}, ids and URLs), the free-text \texttt{description},
which often restates the price, scrape timestamps, and \texttt{views}, which accumulates \emph{after} listing.
\texttt{run\_all} asserts that no leakage column survives into the feature names.

\subsection{Features, tasks and split}
The \NFeat{} features are: numeric attributes (area, rooms, floor, total floors, land, lat/lng); engineered
features (floor ratio, top and ground floor flags, area per room, haversine distance to the city centre); binary amenities
(repair, bill of sale, mortgage, agent, residential complex, VIP/featured); and one-hot category,
district and neighbourhood columns. Neighbourhoods with fewer than 100 training listings are pooled into ``other''.
We chose one-hot over ordinal codes because a linear SVM would otherwise read a fake ordering into nominal values.
The split is 70/15/15 and stratified on price deciles, which balances the tier label across splits (Table~\ref{tab:data}).
The tier threshold is the median of the \emph{training} prices only. Task~A is evaluated with RMSE, MAE and
$R^2$ on log price and on AZN. Task~B is evaluated with precision, recall, F1, ROC-AUC, PR-AUC and the confusion matrix.

\begin{table}[t]
  \centering
  \caption{Dataset summary after cleaning.}
  \label{tab:data}
  \small
  \begin{tabular}{@{}lr@{}}
\toprule
Property & Value \\
\midrule
Raw rows (scraped records) & 100,775 \\
Rows after cleaning & 57,422 \\
Features (after one-hot) & 99 \\
Train / Val / Test & 40,196 / 8,613 / 8,613 \\
Tier threshold (train median) & 209,000 AZN \\
Premium share train / val / test (\%) & 49.8 / 49.9 / 49.8 \\
\bottomrule
\end{tabular}

\end{table}

\begin{figure}[t]
  \centering
  \includegraphics[width=\linewidth]{eda_price_distribution.png}\\[2pt]
  \includegraphics[width=\linewidth]{eda_location.png}
  \caption{Price is heavy-tailed; $\log(\text{price})$ is close to symmetric (top). Location matters (bottom):
  the seaside villa district Khazar and the central Sabail, Nasimi and Narimanov have the highest median
  prices, while the outer Absheron, Qaradagh and Surakhani have the lowest.}
  \label{fig:eda}
\end{figure}

\section{Method}
\subsection{Decision tree (from scratch)}
\texttt{DecisionTree(task, criterion, max\_depth, min\_samples\_split, min\_samples\_leaf,
min\_impurity\_decrease)} mirrors sklearn's \texttt{fit/predict/predict\_proba}. For every node and feature
we sort the samples once. Prefix sums of class counts (classification) or of $y$ and $y^2$ (regression)
then give the impurity of \emph{every} candidate split in one vectorised pass. Candidates lie only between two
distinct values, the threshold is their midpoint, and children smaller than \texttt{min\_samples\_leaf}
are excluded. The growth is recursive and stops at \texttt{max\_depth}, below \texttt{min\_samples\_split},
at a pure node, or when the weighted decrease $\tfrac{N_t}{N}\mathrm{IG}$ falls below
\texttt{min\_impurity\_decrease} (sklearn's definition). A leaf predicts the majority class (with
\texttt{predict\_proba} from leaf class frequencies) or the mean. Feature importances are the normalised
sums of $N_t\cdot\mathrm{IG}$. Two numerical details turned out to matter. First, ties go to the lowest feature index, whereas
sklearn permutes features at random. Second, when two neighbouring values are adjacent doubles, their
midpoint rounds onto the upper value, so we fall back to the lower one. Our real data triggered this bug and it is
now covered by a test.

\subsection{Support-vector machine (from scratch)}
\texttt{PegasosSVM} runs mini-batch Pegasos (batch 32; batch 1 and full batch are also supported). Each
step applies $\mathbf{w}\leftarrow(1-\eta_t\lambda)\mathbf{w}+\tfrac{\eta_t}{|B|}\sum_{i\in A}y_i\mathbf{x}_i$
over the margin violators $A=\{i:y_i\mathbf{w}^{\top}\mathbf{x}_i<1\}$. The bias is folded in as a constant
feature, as the brief suggests. Optional steps are the projection onto $\lVert\mathbf{w}\rVert\le1/\sqrt\lambda$
and averaging of the second-half iterates. The averaging is the variant covered by the Pegasos guarantee, and in practice it removes the
noise of the $1/(\lambda t)$ schedule. The model exposes $\lambda$, the epochs, the batch size and the learning-rate schedule
($1/(\lambda t)$, $\eta_0/\sqrt t$ or constant). It records $J$ after every epoch and reports $\lVert\mathbf{w}\rVert$, the margin $2/\lVert\mathbf{w}\rVert$
and the support vectors ($y_i f(\mathbf{x}_i)\le1$). Features are standardised with training statistics.

\emph{Non-linearity: option (ii), random Fourier features.} We map
$\phi(\mathbf{x})=\sqrt{2/D}\cos(W\mathbf{x}+\mathbf{c})$ with $W_{ij}\sim\mathcal N(0,2\gamma)$ and
$c_j\sim U[0,2\pi)$, so that $\phi(\mathbf{x})^\top\phi(\mathbf{z})\approx e^{-\gamma\lVert\mathbf{x}-\mathbf{z}\rVert^2}$~\cite{rff},
and run the same linear Pegasos on $\phi(X)$ with $D=1000$. We chose RFF over kernelised Pegasos because, with
$\sim$40k training rows, kernel Pegasos needs $O(n^2)$ kernel evaluations, while RFF costs $O(nD)$ and reuses
the tested linear solver unchanged.

\subsection{Tests}
The pytest suite (32 tests) covers the following. Our tree reproduces sklearn's exact split list
(feature, threshold, depth) and \texttt{predict\_proba} on a 12-point example for Gini, entropy and MSE.
Each of the four hyperparameters changes the tree, and feature importances match sklearn's. Pegasos decreases $J$,
agrees with \texttt{SVC} ($C=1/(\lambda n)$) on more than 97\% of a toy split with weight cosine above 0.99, and RFF approximates the
RBF kernel and solves XOR. All metrics equal sklearn's.

\section{Experimental Setup}
All tuning uses the validation split, and the test split is touched once at the end. For the tree we swept
\texttt{max\_depth}$\times$\texttt{min\_samples\_leaf} for Task~A and criterion$\times$\texttt{max\_depth}, then
\texttt{min\_samples\_leaf}, for Task~B. For the SVM we swept $\lambda\in[10^{-6},10^{-1}]$, compared schedules and epochs,
and ran a $\gamma\times\lambda$ grid for RFF. For the SVM we pick the \emph{most regularised} setting whose
validation F1 is within 0.002 of the best. Smaller differences are below the noise level, and very small $\lambda$
had not converged in 30 epochs (Table~\ref{tab:svm}). The final configurations are:
tree A depth~\TreeADepth, leaf~$\ge$\TreeALeaf; tree B \TreeBCrit, depth~\TreeBDepth, leaf~$\ge$\TreeBLeaf;
linear SVM $\lambda=\SvmLam$; RFF $\gamma=\RffGamma$, $\lambda=\RffLam$.
The baselines use the \emph{same} hyperparameters: \texttt{DecisionTreeRegressor/Classifier}, and
\texttt{Pipeline(StandardScaler, SGDClassifier(hinge, alpha=}$\lambda$\texttt{))} as the matched linear SVM.
\texttt{SVC} (linear and RBF, $C=1/(\lambda n)$, same $\gamma$) is fit on a stratified 10k subsample because its
training time grows super-linearly (Fig.~\ref{fig:scaling}). Uncertainty is reported in two ways: a bootstrap std
on the test set (200 resamples), and 5-fold CV mean$\pm$std on train+val, with scaling refit inside each fold.
The imputation medians and one-hot vocabularies are fit once on the training split and reused in the CV folds. This
is a mild optimism for the folds drawn from the validation rows, and it does not affect the test numbers.

\section{Results}
\subsection{Tuning and the bias--variance story}
Fig.~\ref{fig:tree} shows the classic picture. Training error keeps falling with depth, but with
leaves of size 1 the validation RMSE bottoms out at depth 12 and then rises as the tree memorises
individual listings. A minimum leaf size of 5--20 removes most of that variance. Gini and entropy are within
noise of each other (best val F1 \GiniBest{} vs \EntropyBest). The learning curves (Fig.~\ref{fig:lc})
show both failure modes. The tuned deep tree has near-zero training error and a large gap that shrinks as data
is added, which is high variance, so more data still helps. The depth-6 tree and the linear SVM have training error
close to validation error, which is high bias: more data does not help them, and only more capacity would.

\begin{figure}[t]
  \centering
  \includegraphics[width=\linewidth]{tree_taskA_validation_curves.png}\\[2pt]
  \includegraphics[width=\linewidth]{tree_taskB_validation_curves.png}
  \caption{Our tree's validation curves. Top: Task A train vs validation RMSE (left) and the effect
  of \texttt{min\_samples\_leaf} (right). Bottom: Gini vs entropy over depth, and the leaf-size sweep.}
  \label{fig:tree}
\end{figure}

\begin{figure}[t]
  \centering
  \includegraphics[width=\linewidth]{learning_curves.png}
  \caption{Learning curves for Task B: shallow tree, tuned deep tree and linear SVM (train vs validation error).}
  \label{fig:lc}
\end{figure}

For the SVM (Table~\ref{tab:svm}, Fig.~\ref{fig:svm}), shrinking $\lambda$ grows $\lVert\mathbf{w}\rVert$, narrows the
margin and reduces the number of support vectors. Validation F1 saturates at $\lambda\le10^{-3}$. Below that,
30 epochs no longer reach the optimum ($J$ rises again in Table~\ref{tab:svm}) because the
$1/(\lambda t)$ steps are huge early on. At the chosen $\lambda$ the averaged, last-iterate and constant-rate runs all
converge to $J\approx0.258$ within 20 epochs. The full-batch run takes only one step per epoch and is still far from the optimum after 50.
More epochs widen the margin towards its converged value, while F1 barely moves after two epochs. In the 2-D illustration
(Fig.~\ref{fig:margin}) one can see the margin band and the support vectors directly.

\begin{table}[t]
  \centering
  \caption{Linear Pegasos SVM: $\lambda$ sweep (validation).}
  \label{tab:svm}
  \footnotesize
  \setlength{\tabcolsep}{3.5pt}
  \begin{tabular}{@{}lrrrrrr@{}}
\toprule
$\lambda$ & val F1 & val AUC & ‖w‖ & 2/‖w‖ & SV \% & J \\
\midrule
1e-01 & 0.871 & 0.954 & 1.16 & 1.731 & 46.0 & 0.393 \\
1e-02 & 0.899 & 0.964 & 2.27 & 0.883 & 31.3 & 0.288 \\
1e-03 & 0.902 & 0.965 & 3.39 & 0.591 & 25.7 & 0.258 \\
1e-04 & 0.903 & 0.965 & 5.91 & 0.338 & 19.8 & 0.265 \\
1e-05 & 0.903 & 0.964 & 25.32 & 0.079 & 11.2 & 0.627 \\
1e-06 & 0.903 & 0.964 & 221.83 & 0.009 & 9.4 & 4.710 \\
\bottomrule
\end{tabular}

\end{table}

\begin{figure}[t]
  \centering
  \includegraphics[width=\linewidth]{svm_lambda_sweep.png}\\[2pt]
  \includegraphics[width=\linewidth]{svm_convergence.png}
  \caption{SVM: $\lambda$ trades margin width against fit (top). The objective $J$ decreases and
  converges, compared across schedules and epochs (bottom).}
  \label{fig:svm}
\end{figure}

\begin{figure}[t]
  \centering
  \includegraphics[width=\linewidth]{svm_margin_2d.png}
  \caption{A margin you can see: our SVM on two standardised features. Solid line: $f=0$; dashed: $f=\pm1$;
  ringed points: support vectors.}
  \label{fig:margin}
\end{figure}

\subsection{Benchmark against scikit-learn}
Tables~\ref{tab:resA}--\ref{tab:resB} and the combined Table~\ref{tab:all} hold the headline results.
On Task~A our tree and sklearn's are within one bootstrap std (\OursARmse{} vs \SkARmse{} RMSE on log price),
and both are far better than predicting the mean (\BaseARmse). In AZN the MAE is about \OursAMae k~AZN.
On Task~B every model lands in a narrow band: our tree \OursBTreeFone{} vs sklearn \SkBTreeFone; our
linear SVM \OursSvmFone{} vs SGD \SgdFone{} and SVC \SvcLinFone; our RFF-SVM \OursRffFone{}
vs RBF-SVC \SvcRbfFone. Our final linear SVM has $\lVert\mathbf{w}\rVert=\SvmNormW$, margin
$2/\lVert\mathbf{w}\rVert=\SvmMargin$, and \SvmSvFrac\% of the training points are support vectors.

\emph{Structure (Table~\ref{tab:struct}).} With matched hyperparameters the two trees choose the same root
split (area $\le$ 100.25~m\textsuperscript{2} for A) and share 6--7 of the first 7 splits. They agree on \SameA\% (A) and \SameB\% (B) of test predictions. The
remaining gaps come from three sources: (a) sklearn stores $X$ as float32, so thresholds are rounded and
nearly-equal values can merge; (b) ties are broken by a random feature permutation in sklearn and by the lowest
index in ours; (c) after an early disagreement the subtrees diverge.

\emph{Speed (Fig.~\ref{fig:scaling}).} Our tree is written in Python over NumPy and fits in \OursTreeAFit~s
vs \SkTreeAFit~s for sklearn's Cython (Task~A), 25--45$\times$ slower across our runs, because every node pays Python overhead
on top of the sort. Both trees grow slightly faster than linearly in $n$. Wall-clock times on a shared laptop are noisy.
Our Pegasos fits in \OursSvmFit~s and is linear in $n$. Kernel SVC is already slower on a 10k subsample (\SvcRbfFit~s) and its
cost grows super-linearly, which is why the RFF route is the practical way to get a kernel at this size.

\begin{table*}[t]
  \centering
  \caption{Task A (log price) on the test set. RMSE $\pm$ bootstrap std; CV = 5-fold mean $\pm$ std on train+val.}
  \label{tab:resA}
  \small
  \begin{tabular}{@{}lrrrrrrrr@{}}
\toprule
Model & RMSE log & MAE log & R$^2$ log & RMSE AZN & MAE AZN & R$^2$ AZN & Fold RMSE log* & fit s \\
\midrule
Baseline: predict train mean & 0.636 $\pm$ 0.006 & 0.490 & -0.000 & 264.7k & 130.1k & -0.046 & n/a & 0.00 \\
Ours: DecisionTree (MSE) & 0.231 $\pm$ 0.004 & 0.142 & 0.869 & 134.7k & 46.1k & 0.729 & 0.234 $\pm$ 0.003 & 17.31 \\
sklearn DecisionTreeRegressor & 0.232 $\pm$ 0.004 & 0.143 & 0.867 & 135.0k & 46.3k & 0.728 & 0.233 $\pm$ 0.002 & 0.70 \\
\bottomrule
\end{tabular}

\end{table*}

\begin{table*}[t]
  \centering
  \caption{Task B (price tier) on the test set. F1 $\pm$ bootstrap std; CV F1 = 5-fold mean $\pm$ std on train+val.
  SVC is fit on a stratified 10k subsample of the training split.}
  \label{tab:resB}
  \small
  \begin{tabular}{@{}lrrrrrrrr@{}}
\toprule
Model & Prec. & Rec. & F1 (test) & ROC-AUC & PR-AUC & Fold F1* & fit s & pred. s \\
\midrule
Ours: DecisionTree (entropy) & 0.922 & 0.924 & 0.923 $\pm$ 0.003 & 0.926 & 0.894 & 0.923 $\pm$ 0.003 & 9.83 & 0.031 \\
sklearn DecisionTreeClassifier & 0.927 & 0.927 & 0.927 $\pm$ 0.003 & 0.930 & 0.900 & 0.923 $\pm$ 0.003 & 1.31 & 0.006 \\
Ours: Pegasos linear SVM & 0.910 & 0.908 & 0.909 $\pm$ 0.003 & 0.968 & 0.965 & 0.906 $\pm$ 0.003 & 4.98 & 0.201 \\
Ours: RFF-Pegasos SVM & 0.912 & 0.913 & 0.913 $\pm$ 0.003 & 0.974 & 0.973 & 0.909 $\pm$ 0.002 & 9.10 & 0.316 \\
sklearn SGDClassifier (hinge) & 0.906 & 0.909 & 0.907 $\pm$ 0.003 & 0.967 & 0.964 & 0.905 $\pm$ 0.003 & 1.83 & 0.010 \\
sklearn SVC linear (10k subsample) & 0.906 & 0.903 & 0.904 $\pm$ 0.003 & 0.966 & 0.963 & 0.905 $\pm$ 0.003 & 6.29 & 2.994 \\
sklearn SVC RBF (10k subsample) & 0.911 & 0.908 & 0.910 $\pm$ 0.003 & 0.973 & 0.973 & 0.908 $\pm$ 0.003 & 7.01 & 4.800 \\
\bottomrule
\end{tabular}

\end{table*}

\begin{table}[t]
  \centering
  \caption{All models on both tasks (test set; CV on train+val).}
  \label{tab:all}
  \scriptsize
  \setlength{\tabcolsep}{2pt}
  \resizebox{\linewidth}{!}{\begin{tabular}{@{}lrrrrrr@{}}
\toprule
Model & A: RMSE log & A: R$^2$ log & A: CV RMSE & B: F1 & B: AUC & B: CV F1 \\
\midrule
Baseline: predict train mean & 0.636 & -0.000 & n/a & — & — & — \\
Ours: DecisionTree (MSE) & 0.231 & 0.869 & 0.234 $\pm$ 0.003 & — & — & — \\
sklearn DecisionTreeRegressor & 0.232 & 0.867 & 0.233 $\pm$ 0.002 & — & — & — \\
Ours: DecisionTree (entropy) & — & — & — & 0.923 & 0.926 & 0.923 $\pm$ 0.003 \\
sklearn DecisionTreeClassifier & — & — & — & 0.927 & 0.930 & 0.923 $\pm$ 0.003 \\
Ours: Pegasos linear SVM & — & — & — & 0.909 & 0.968 & 0.906 $\pm$ 0.003 \\
Ours: RFF-Pegasos SVM & — & — & — & 0.913 & 0.974 & 0.909 $\pm$ 0.002 \\
sklearn SGDClassifier (hinge) & — & — & — & 0.907 & 0.967 & 0.905 $\pm$ 0.003 \\
sklearn SVC linear (10k subsample) & — & — & — & 0.904 & 0.966 & 0.905 $\pm$ 0.003 \\
sklearn SVC RBF (10k subsample) & — & — & — & 0.910 & 0.973 & 0.908 $\pm$ 0.003 \\
\bottomrule
\end{tabular}
}
\end{table}

\begin{table}[t]
  \centering
  \caption{Our tree vs sklearn's with matched hyperparameters.}
  \label{tab:struct}
  \footnotesize
  \setlength{\tabcolsep}{3pt}
  \resizebox{\linewidth}{!}{\begin{tabular}{@{}lrrrr@{}}
\toprule
 & depth (ours/sk) & leaves (ours/sk) & first 7 splits equal & same test pred. \\
\midrule
Task A & 20 / 20 & 5324 / 5337 & 6/7 & 93.6\% \\
Task B & 20 / 20 & 2231 / 2239 & 7/7 & 98.2\% \\
\bottomrule
\end{tabular}
}
\end{table}

\begin{figure}[t]
  \centering
  \includegraphics[width=0.85\linewidth]{scaling_fit_time.png}
  \caption{Fit time vs training size (log--log).}
  \label{fig:scaling}
\end{figure}

\subsection{Error analysis}
\emph{What the models use.} The tree's impurity-decrease importances (Fig.~\ref{fig:imp}) are dominated by
area (0.59) and location (distance to the centre 0.23, then longitude and latitude), followed by land area
and category. Our importances and sklearn's are almost identical. The SVM weights tell the same story with
signs: area, rooms and area per room push towards ``premium'', while distance from the centre and outer settlements
such as Bin\schwa{} push towards ``standard''.

\emph{Where the models fail.} (i) \emph{Luxury and bargain tails:} the tree shrinks towards the middle. It
over-predicts the cheapest decile and under-predicts the most expensive one (Fig.~\ref{fig:errreg}, left),
because a leaf mean cannot extrapolate. The largest AZN errors (Table~\ref{tab:worst}) are 700--1200~m\textsuperscript{2}
villas listed at 3--5M~AZN, plus one large new-build flat. (ii) \emph{Houses and the districts dominated by them:} the MAE on log price for houses is about 2.6$\times$ that
for flats (Fig.~\ref{fig:errreg}, right), because land, condition and plot shape vary far more and are only
partly captured by our features. Per district (12 districts with at least 10 test rows), the test MAE correlates
strongly with the share of houses ($r=\DistHouseR$) and more weakly with the number of training listings
($r=\DistSizeR$). The smaller districts are also the suburban, house-heavy ones, so the two effects overlap.
House share is the better explanation: Khazar, almost entirely villas, is the worst district despite 2.7k
training listings. With only 12 points this is suggestive, not conclusive. (iii) Flats in the dense
central districts have the smallest errors (MAE $\approx$0.09--0.12 on log price, i.e.\ roughly 10\%).
(iv) \emph{Tier errors near the threshold:} \SvmErrNear\% of the SVM's tier errors are listings within about 20\% of
the \TierThr~AZN threshold (Fig.~\ref{fig:errcls}). Many of these are irreducible, because asking prices of near-identical
flats differ by more than that. Near the threshold the deep tree errs less than both SVMs, since it can carve
out local neighbourhoods. Far from the threshold the smoother SVMs make fewer mistakes than the tree, whose leaves sometimes
hold a few mislabelled outliers. This is consistent with the SVMs' higher ROC-AUC despite a lower F1.

\begin{figure}[t]
  \centering
  \includegraphics[width=\linewidth]{importance_and_weights.png}
  \caption{Left: impurity-decrease importances (Task A), ours vs sklearn. Right: largest linear-SVM
  weights (Task B; positive pushes towards premium).}
  \label{fig:imp}
\end{figure}

\begin{figure}[t]
  \centering
  \includegraphics[width=\linewidth]{errors_regression.png}
  \caption{Task A test errors of our tree: mean residual by true-price decile (left), district MAE vs the
  share of houses in the district (middle), and MAE by category (right).}
  \label{fig:errreg}
\end{figure}

\begin{figure}[t]
  \centering
  \includegraphics[width=\linewidth]{roc_confusion.png}\\[2pt]
  \includegraphics[width=0.8\linewidth]{errors_classification_threshold.png}
  \caption{Task B: ROC curves and confusion matrices (test). Error rate vs distance of the true
  price from the tier threshold.}
  \label{fig:errcls}
\end{figure}

\begin{table}[t]
  \centering
  \caption{Largest absolute AZN errors of our Task A tree (test).}
  \label{tab:worst}
  \footnotesize
  \begin{tabular}{@{}llrrr@{}}
\toprule
category & district & m$^2$ & true & pred. \\
\midrule
house & Nerimanov r. & 1000 & 5,000k & 1,562k \\
house & Sabuncu r. & 1200 & 4,500k & 1,301k \\
house & Sebail r. & 1000 & 3,500k & 1,214k \\
house & Xezer r. & 1000 & 4,000k & 2,105k \\
house & Xezer r. & 1000 & 4,000k & 2,105k \\
house & Xezer r. & 1000 & 4,000k & 2,105k \\
new build & Nesimi r. & 237 & 3,094k & 1,232k \\
house & Sebail r. & 700 & 3,000k & 1,214k \\
\bottomrule
\end{tabular}

\end{table}

\section{Discussion and Limitations}
\begin{itemize}\setlength\itemsep{0pt}
\item \emph{Asking, not transaction, prices.} The target is the listed price. Agents' over-pricing and stale
  listings add noise that no feature can explain.
\item \emph{One snapshot.} The data covers five weeks in autumn 2024. A random split mixes the same buildings
  into train and test, and a time- or building-grouped split would be harder. We did not model time.
\item \emph{Unused text.} \texttt{description} probably holds strong signal (renovation quality, view), but it
  sometimes restates the price, so we excluded it rather than risk leakage.
\item \emph{Speed.} Our tree is 25--45$\times$ slower than sklearn's Cython. This is acceptable here, but a
  histogram-based split search would be needed for much larger data.
\item \emph{SVC on a subsample.} The kernel-SVC baseline sees 10k rows, not 40k. That favours our models slightly,
  and we state it in every table.
\item \emph{Residual duplicates.} Three identical rows in Table~\ref{tab:worst} (same district, area and price)
  suggest that re-posts with slightly different coordinates survive our exact-match de-duplication. A fuzzy
  match (e.g.\ same price and area within 200~m) would be stricter.
\item \emph{Selection rule.} Choosing the most regularised SVM within 0.002 F1 is a judgement call. The
  validation-best $\lambda=10^{-6}$ would give the same F1 within noise, but with an unconverged objective.
\end{itemize}

\section{Conclusion}
From-scratch implementations of CART and Pegasos, checked against sklearn on toy and real data, reach the
same accuracy as the library versions on a real, messy dataset once leakage and duplicates are handled. Location
and area explain most of the price. The remaining error sits in the tails and close to the tier threshold.
Next steps are a grouped/temporal split, text features from descriptions with price mentions masked,
and ensembles of our own trees.

\section*{Tools and Acknowledgements}
\emph{AI assistance:} an AI coding assistant (Claude, Anthropic) was used to draft the code, tests,
this report and the slides. All results were produced by running the code in the repository, and the team
reviewed and can explain every component (see \texttt{contribution\_report.md}).
\emph{Data:} bina.az listings via Kaggle (\texttt{sehriyarmemmedli/binaaz-sale-project}). Personal fields were
dropped and are not redistributed. \emph{Libraries:} NumPy, pandas and Matplotlib, with scikit-learn used for baselines and tests only.

\begin{thebibliography}{9}
\bibitem{pegasos}
S.~Shalev-Shwartz, Y.~Singer, N.~Srebro, and A.~Cotter,
``Pegasos: primal estimated sub-gradient solver for SVM,''
\emph{Mathematical Programming}, vol.~127, no.~1, pp.~3--30, 2011.
\bibitem{rff}
A.~Rahimi and B.~Recht, ``Random features for large-scale kernel machines,''
in \emph{Proc. NeurIPS}, 2007, pp.~1177--1184.
\bibitem{cart}
L.~Breiman, J.~Friedman, R.~Olshen, and C.~Stone, \emph{Classification and Regression Trees}.
Wadsworth, 1984.
\bibitem{sklearn}
F.~Pedregosa \emph{et al.}, ``Scikit-learn: machine learning in Python,''
\emph{JMLR}, vol.~12, pp.~2825--2830, 2011.
\end{thebibliography}

\end{document}
